\section{Renormalisation freedom and invariant content}
\label{sec:qs-freedom}

\begin{definition}[The arity-one finite counterterm along $\Sigma$]
\label{def:qs-Fxi}
For $\xi\in\Sigma$ write the quadratic part of $S^{(1)}(\xi)$ as $S^V_2(\xi)+S^\perp_2(\xi)$, with $S^V_2(\xi)=T_1+T_3$ the part with the legs $(\lambda,A^*|_V)$, $(\AV,\lambda^*)$ and $S^\perp_2(\xi)=T_2+T_4$ the part with the legs $(\lambda,A^*|_\perp)$, $(\Ap,\lambda^*)$ (Table \ref{tab:qs-terms}). The arity-one equation along $\xi$ is
\begin{equation}
(Q+\{I,-\})\,\hbar F_1^{(1)}(\xi)+\big\{S^{(1)}(\xi),\hbar F_1^{(0)}\big\}=\hbar\,\mathrm{Ob}^{(1)}(\xi),
\label{eq:qs-arity1xi}
\end{equation}
the arity-one part along $\xi$ of $s(\hbar F_1)=\hbar\,\mathrm{Ob}$ (Thm c8-obstruction (ii)); $\hbar\,\mathrm{Ob}^{(1)}(\xi)$ is the local functional whose tree flow is the order-$\hbar$ limit of the arity-one defect of the action along $\xi$: the cycles of vertices of $I$ with one vertex $S^{(1)}(\xi)$ and one special edge (the classes (B), (D) without $O$), and the counterterm vertices $-QI^{CT}_{\mathfrak{psl}(4|4)}[\varepsilon]^{(1)}(\xi)-\{I,I^{CT}_{\mathfrak{psl}(4|4)}[\varepsilon]^{(1)}(\xi)\}_\varepsilon-\{S^{(1)}(\xi),I^{CT}[\varepsilon]\}_\varepsilon$ (Prop o-uv).
\end{definition}

\begin{proposition}[The quadratic free directions along $\Sigma$]
\label{prop:qs-freexi}
\begin{itemize}
\item (i) The quadratic, $\mathrm{Spin}(4)\times\mathrm{Spin}(6)\times\mathfrak g$-covariant local functionals of degree $-1$ and Weyl weight $\tfrac92$ (fields and derivatives), linear in $\xi\in S_+\otimes W$ and translation invariant, are spanned by the five functionals $T_1$, $T_2$, $T_3$, $T_4$ and $\int\big((d^*A|_V)\,\xi,\lambda^*\big)$; their $Q$-closed combinations are spanned by $S^V_2(\xi)$ and $S^\perp_2(\xi)$, and no nonzero combination is $Q$-exact.
\item (ii) Two solutions of \eqref{eq:qs-arity1xi} differ by the arity-one part along $\xi$ of an $s$-closed cochain of degree $0$; in the quadratic sector by $\theta^\xi_VS^V_2(\xi)+\theta^\xi_\perp S^\perp_2(\xi)$. The cubic part of \eqref{eq:qs-arity1xi} contains $\theta^\xi_V\{I_3,S^V_2(\xi)\}+\theta^\xi_\perp\{I_3,S^\perp_2(\xi)\}$, and since $\{I,S^{(1)}(\xi)\}=-QS^{(1)}(\xi)$ only the combination with $\theta^\xi_V=\theta^\xi_\perp$ extends, with the cubic term $\theta^\xi S_3(\xi)$, to an $s$-closed direction of the cochain: $\theta^\xi S^{(1)}(\xi)$, the arity-one value of $G_B$ (Def c8-B) and of the field redefinition by the R-rotation $S^{(1)}(R_\xi)$, $[R_\xi,\xi]=\xi$ (Thm c8-twist).
\item (iii) The one-body kernel of a letter $X$ along $\xi$ is $-\{\hbar F_1^{(1)}(\xi),X\}$; the free directions add $-\theta^\xi_V\{S^V_2(\xi),X\}-\theta^\xi_\perp\{S^\perp_2(\xi),X\}$, which on the letters $\AV$, $\lambda$ (the $V$-arrow), $\Ap$, $\lambda$ (the $\perp$-arrow) is $\theta^\xi_V$ or $\theta^\xi_\perp$ times the corresponding arrow of the classical transformation.
\end{itemize}
The same holds along $S_-\otimes W^*$ with the chiralities exchanged, along $\Sigma^*$ for the position part with $\rho(b)\eta$ in place of $\xi$, and along $\Sigma^*$ for the intrinsic part with the additional direction $\int s_0$.
\end{proposition}

\begin{proof}
(i) Field contents of degree $-1$ and field weight $\tfrac92$ linear in a constant spinor: $(\lambda,A^*)$ without derivatives and $(A,\lambda^*)$ with one derivative ($c$ and $c^*$ would need weight $\tfrac12$); covariant pairings: $\Gamma(\xi,\lambda)$ with $A^*$ in its $V$ and $W^\perp_{\mathbb R}$ parts, and $\rho(\partial A)\xi$ with $\lambda^*$, where $\partial A$ is $dA|_V$, $d^*A|_V$ or the mixed part $\rho(e_m)\rho(\partial_mA|_\perp)$; a factor $\rho(e_1e_2e_3e_4)$ acts on each chirality summand by a sign and gives no new structure. On functionals $Q$ substitutes $A\mapsto-dc$, $A^*\mapsto-\tfrac12d^*dA$, $\lambda^*\mapsto{\not d}\lambda$, $c^*\mapsto d^*(A^*|_V)$ and annihilates $c$, $\lambda$ (Prop c1-Q); $Q$ of $\int((d^*A|_V)\xi,\lambda^*)$ contains $-\int((d^*dc)\xi,\lambda^*)$, produced by no other structure, so its coefficient vanishes in a closed combination; the remaining four map to $(\lambda,A)$-functionals, and $Q$ preserves the split into the parts along $V$ and along $W^\perp_{\mathbb R}$, so the closed space is at most two-dimensional; $S^V_2(\xi)$ and $S^\perp_2(\xi)$ are closed, being the two parts of the quadratic part of $\{S,S^{(1)}(\xi)\}=0$. A $Q$-exact functional would be $Q$ of a functional of degree $-2$ and field weight $\tfrac92$, which does not exist (two antifields have weight at least $5$, and $c^*$ needs a further field of weight $\tfrac12$). (ii) Thm c8-obstruction (iii) restricted to arity one along $\xi$; the cubic part of $s(\theta^\xi_VS^V_2+\theta^\xi_\perp S^\perp_2)$ is the bracket with $I_3$, and $\{I_3,S_2(\xi)\}=-QS_3(\xi)$ for the full quadratic part only; $G_B^{(1)}(\xi)=S^{(1)}(B\xi)=S^{(1)}(\xi)$ for $\xi\in S_+\otimes W$ (Def c8-B). (iii) Theorem \ref{thm:qs-classes} (iii) and Lemma \ref{lem:qs-arrows}. The engine task TC0-$\Sigma$ of Section \ref{sec:qs-engines} confirms (i) by linear algebra on plane-wave normal forms, as Prop c11-tc0 does along $\Delta$.
\end{proof}

\begin{proposition}[Transformation laws of $Q_1$, $S_1$, $K_1$]
\label{prop:qs-laws}
Under the changes of Prop c8-laws the components change as follows; $\{S^{(1)}(x),-\}$ is the classical transformation.
\begin{itemize}
\item (a) Field redefinition $F\mapsto F+\hbar s\phi$: the component along $x$ changes by $-\{(s\phi)^{(1)}(x),-\}$, $(s\phi)^{(1)}(x)=\pm\{S^{(1)}(x),\phi^{(0)}\}+(Q+\{I,-\})\phi^{(1)}(x)$: a one-body change. For the field countings $\phi^{(0)}=\sum_fy_f\int(f,f^*)$ the change of $Q_1$ along $\xi$ multiplies each arrow $f\to f'$ of the classical transformation by $y_{f'}-y_f$ (Prop c8-B proof); for $\phi^{(0)}=t\,S^{(1)}(\Delta)$, $(s\phi)^{(1)}(x)=t\{S^{(1)}(x),S^{(1)}(\Delta)\}$ is $\pm t\lambda_x S^{(1)}(x)$, the sign fixed by $\{S^{(1)}(\xi),S^{(1)}(\Delta)\}=S^{(1)}(\xi)$ and $\{S^{(1)}(P_b),S^{(1)}(\Delta)\}=-2S^{(1)}(P_b)$ (proof of Prop c8-B), so the change is a multiple of the classical transformation proportional to $\lambda_x$, zero along $\Delta$ and the rotations.
\item (b) Finite parts of the operators $\psi\mapsto\psi+\psi'$: the component along $x$ changes by $[\{S^{(1)}(x),-\},\psi']$, i.e. on $O$ by $\{S^{(1)}(x),\psi'(O)\}-\psi'(\{S^{(1)}(x),O\})$. For $\psi'=\tfrac a2\{S^{(1)}(\Delta),-\}$ this is $\tfrac a2\lambda_x\{S^{(1)}(x),-\}$: $-\tfrac a2$ along $\Sigma$, $+\tfrac a2$ along $\Sigma^*$, $+a$ along $K_b$ (Prop c9-scheme), a one-body change. A finite part $\psi'$ of the operators of the Weyl weight of $\{S^{(1)}(x),O\}$ that vanishes on $O$ changes the component by $-\psi'(\{S^{(1)}(x),O\})$, a change of the two-body kernels: for the Konishi descendant $O_\zeta=\int\omega\,(u'^\alpha\lambda_\alpha{}^i,A_{ij}z^j)$, whose classical component along $\xi=u\otimes z$ is $\kappa_B$ times the cubic $B_z=\int\omega\,([A^{ik},A_{kj}z^j],A_{il}z^l)$, the finite part $\psi'(B_z)=y\,\bar F_z$, $\bar F_z=\int\omega\,(\bar\lambda_{\dot\alpha i}z^i,\bar\lambda^{\dot\alpha}{}_jz^j)$, changes the Konishi entry $c_K$ by $-\kappa_By$ (Rem c11-twist, Prop o-konishi).
\item (c) Closed directions: the coupling direction leaves $Q_1$, $S_1$, $K_1$ unchanged and changes $L_1^{(0)}$ by $-a(Q+\{I,-\})$ and $L_1^{(2)}$ by $a\{S^{(2)},-\}$; the outer direction $F\mapsto F+\hbar aG_B$ changes $Q_1$ along $S_+\otimes W$ by $-a\{S^{(1)}(\xi),-\}$ and along $S_-\otimes W^*$ by $+a\{S^{(1)}(\xi),-\}$, $S_1$ along $S_+\otimes W^*$ by $+a\{S^{(1)}(\eta),-\}$ and along $S_-\otimes W$ by $-a\{S^{(1)}(\eta),-\}$, and leaves $K_1$, $D_1$, $P_1$, $M_1$, $R_1$ unchanged (Prop c8-laws (c)).
\item (d) The free directions of $\hbar F_1^{(1)}(x)$: along $\Sigma$ the one-body changes of Proposition \ref{prop:qs-freexi} (iii); along $\Delta$, $P_b$ the sector coefficients $\theta,\theta_\lambda,\theta_V,\theta_P,\dots$ of Def c11-free, which enter $D_1$, $P_1$, and through $P_1$ the relations of $S_1$ to $Q_1$ (Prop c11-S (ii)) and of $K_1$ to $D_1$ (Proposition \ref{prop:qs-K} (iii)); along $K_b$, $s$-closed directions of degree $-1$ quadratic in the Killing data, which change $K_1$ by one-body terms.
\item (e) The cutoff shape (Prop c10-shape) changes the finite parts of the defect cycles, of the special-edge cycles of $L_1^{(0)}$ and the finite special vertices; by the uniqueness of the counterterms up to finite local terms [C07, Thm B], a change of shape is a change of $\hbar F_1$ and of $\psi$, hence a combination of (a)--(d).
\end{itemize}
\end{proposition}

\begin{proof}
(a)--(c): Prop c8-laws, Prop c9-scheme and Prop c11-freedom, with $[\{S^{(1)}(x),-\},\psi'](O)=\{S^{(1)}(x),\psi'(O)\}-\psi'(\{S^{(1)}(x),O\})$; for $\phi^{(0)}=tS^{(1)}(\Delta)$, $\{S^{(1)}(x),S^{(1)}(\Delta)\}$ is $S^{(1)}$ of $\pm[x,\Delta]=\mp\lambda_xx$ by Lemma c2-components with $S^{(2)}(x,\Delta)=0$ (Def c2-S2). The Konishi entry: Prop c11-freedom (iii) with $\psi'(B_z)=y\bar F_z$, the two functionals having Weyl weight $3$, degree $0$ and the same quantum numbers. (d) Definition \ref{def:qs-Fxi}, Proposition \ref{prop:qs-freexi}, Def c11-free, Prop c11-S (ii), Proposition \ref{prop:qs-K}. (e) Prop c10-shape: only the finite parts change; the finite counterterm and the finite parts of the operators absorb the change, the counterterms being unique up to finite local terms.
\end{proof}

\begin{theorem}[Invariant content]
\label{thm:qs-invariant}
\begin{itemize}
\item (i) Under (a), (b) and the coupling direction the representation of $\mathfrak{psl}(4|4)$ induced on the cohomology of $\Lambda^{(0)}$ by $-\Lambda^{(1)}$ changes by an isomorphism; under $G_B$ it is composed with $e^{-a\hbar B}$ (Thm c8-invariants). In particular the one-loop components along $\Sigma$, $\Sigma^*$ and $K_b$ are determined on cohomology up to conjugation, up to the rescaling $1\mp a\hbar$ of the odd generators by the outer direction, and, through $P_1$, up to the relations of (d).
\item (ii) For a rank-one $\xi=u\otimes z$ the class $[Q^\xi_1]$ on the cohomology of $Q^\xi_0$, its kernel, image and rank are invariant under (a), (b), the coupling direction and $G_B$ (Thm c8-twist); the individual two-body kernels $\Upsilon_\xi$ and the one-body kernels are not.
\item (iii) The relations of Cor c8-homotopy among $Q_1$, $S_1$, $D_1$, $K_1$, $P_1$, $M_1$, $R_1$ on cohomology, and the relations $S_1=-Q_1^{\rm adj}+(\text{Gram terms})$, $K_1=-P_1^{\rm adj}+(\text{Gram terms})$ of Cor c8-gram in a reflection-compatible scheme, are covariant under all the changes of Proposition \ref{prop:qs-laws} with reflection-compatible parameters (real $a$ in (b), $\bar a=-a$ for $G_B$, Cor c8-gram).
\item (iv) The spectrum of $D_1$ on cohomology in a fixed classical weight is invariant (Thm c8-invariants (iii)); with the relations of (iii) it fixes the parts of $S_1$ and $K_1$ that multiply the position on cohomology (Prop c11-S (iii), Proposition \ref{prop:qs-K} (iii)).
\end{itemize}
\end{theorem}

\begin{proof}
Thm c8-invariants, Thm c8-twist, Cor c8-homotopy, Cor c8-gram and Proposition \ref{prop:qs-laws}; each change acts on $\Lambda$ by a conjugation, by a multiple of $\Lambda^{(0)}$ or by the automorphism $e^{-a\hbar B}$, and on the inner product by the corresponding change of $\mathrm{Gram}$ (Cor c8-gram). (iv): Prop c11-S (iii) and Proposition \ref{prop:qs-K} (iii).
\end{proof}
